% !TEX root = ../main.tex
\section{Ejercicio propuesto 1: implementación y validación del montículo de Fibonacci}

\subsection{Implementación de la estructura de datos}
Se implementaron las clases \texttt{FibonacciNode} y \texttt{FibonacciHeap} en Python con punteros circulares bidireccionales, relaciones padre-hijo, grado de ramificación, banderas de corte y validación de ocho invariantes estructurales.

\begin{lstlisting}[language=Python, caption={Implementación compacta del Montículo de Fibonacci en Python}, label={lst:fib_heap_propuesto}]
class FibonacciNode:
    def __init__(self, key, value=None):
        self.key, self.value = key, value
        self.parent = self.child = None
        self.left = self.right = self
        self.degree, self.mark = 0, False

class FibonacciHeap:
    def __init__(self):
        self.min_node = None
        self.n = self.links = self.cuts = self.cascading_cuts = self.consolidations = 0

    def is_empty(self): return self.min_node is None
    def minimum(self): return self.min_node

    def insert(self, key, value=None):
        node = FibonacciNode(key, value)
        if self.min_node is None: self.min_node = node
        else:
            self._insert_root(node)
            if node.key < self.min_node.key: self.min_node = node
        self.n += 1
        return node

    def _insert_root(self, node):
        node.parent = None
        node.left = self.min_node; node.right = self.min_node.right
        self.min_node.right.left = node; self.min_node.right = node

    def union(self, other: 'FibonacciHeap'):
        if other is None or other.min_node is None: return
        if self.min_node is None:
            self.min_node, self.n = other.min_node, other.n
        else:
            s_next, o_prev = self.min_node.right, other.min_node.left
            self.min_node.right = other.min_node; other.min_node.left = self.min_node
            s_next.left = o_prev; o_prev.right = s_next
            if other.min_node.key < self.min_node.key: self.min_node = other.min_node
            self.n += other.n
        other.min_node, other.n = None, 0

    def extract_min(self):
        z = self.min_node
        if z is not None:
            if z.child is not None:
                for x in self._get_circular(z.child):
                    x.parent = None; x.mark = False
                    self._insert_root(x)
                z.child = None
            if z.right == z: self.min_node = None
            else:
                z.left.right, z.right.left = z.right, z.left
                self.min_node = z.right
                self._consolidate()
            self.n -= 1
        return z

    def _get_circular(self, start):
        if start is None: return []
        nodes, curr = [], start
        while True:
            nodes.append(curr); curr = curr.right
            if curr == start: break
        return nodes

    def _consolidate(self):
        self.consolidations += 1
        A = {}
        for w in self._get_circular(self.min_node):
            x, d = w, w.degree
            while d in A:
                y = A[d]
                if x.key > y.key: x, y = y, x
                self._link(y, x); del A[d]; d += 1
            A[d] = x
        self.min_node = None
        for node in A.values():
            node.left = node.right = node; node.parent = None
            if self.min_node is None: self.min_node = node
            else:
                self._insert_root(node)
                if node.key < self.min_node.key: self.min_node = node

    def _link(self, y, x):
        self.links += 1
        y.left.right, y.right.left = y.right, y.left
        y.parent = x
        if x.child is None: x.child = y.left = y.right = y
        else:
            y.left = x.child; y.right = x.child.right
            x.child.right.left = y; x.child.right = y
        x.degree += 1; y.mark = False

    def decrease_key(self, x, k):
        if k > x.key: raise ValueError(f"Nueva clave {k} > {x.key}")
        if k == x.key: return
        x.key = k; y = x.parent
        if y is not None and x.key < y.key:
            self._cut(x, y); self._cascading_cut(y)
        if self.min_node is None or x.key < self.min_node.key: self.min_node = x

    def _cut(self, x, y):
        self.cuts += 1
        if x.right == x: y.child = None
        else:
            if y.child == x: y.child = x.right
            x.left.right, x.right.left = x.right, x.left
        y.degree -= 1; self._insert_root(x)
        x.parent, x.mark = None, False

    def _cascading_cut(self, y):
        z = y.parent
        if z is not None:
            if not y.mark: y.mark = True
            else:
                self.cascading_cuts += 1
                self._cut(y, z); self._cascading_cut(z)

    def delete(self, x):
        self.decrease_key(x, float('-inf'))
        self.min_node = x
        return self.extract_min()

    def validate(self, after_extract_min=False):
        if self.min_node is None: assert self.n == 0; return True
        roots = self._get_circular(self.min_node)
        assert self.min_node.key == min(r.key for r in roots)
        for r in roots: assert not r.mark and r.parent is None
        if after_extract_min:
            degs = [r.degree for r in roots]
            assert len(degs) == len(set(degs))
        visited = set()
        for r in roots: self._validate_sub(r, visited)
        assert len(visited) == self.n
        return True

    def _validate_sub(self, node, visited):
        assert id(node) not in visited; visited.add(id(node))
        assert node.right.left == node and node.left.right == node
        if node.parent: assert node.parent.key <= node.key
        ch = self._get_circular(node.child)
        assert len(ch) == node.degree
        for c in ch: assert c.parent == node; self._validate_sub(c, visited)
\end{lstlisting}

\subsection{Batería de pruebas y validación de invariantes}
Se ejecutó la totalidad de casos de prueba de la guía, evaluando casos frontera, ordenación con repetidos, fusión de montículos, cortes en cascada multinivel, eliminación interna y 250 operaciones aleatorias.

\begin{lstlisting}[language=Python, caption={Código de la batería de pruebas y validaciones obligatorias}, label={lst:pruebas_propuesto_1}]
# 1. Extraccion vacio y nodo unico
h = FibonacciHeap(); assert h.is_empty() and h.extract_min() is None
h.insert(99); assert h.extract_min().key == 99 and h.is_empty()

# 2. Inserciones con duplicados y ordenacion
claves = [10, 20, 30, 5, 4, 3, 15, 25, 20, 10, 5]
for k in claves: h.insert(k)
ext = [h.extract_min().key for _ in range(len(claves))]
assert ext == sorted(claves)

# 3. Union de monticulos
h1, h2 = FibonacciHeap(), FibonacciHeap()
for k in [50, 10, 30]: h1.insert(k)
for k in [40, 20, 60]: h2.insert(k)
h1.union(h2); assert h1.n == 6 and h1.minimum().key == 10 and h2.is_empty()

# 4. Disminucion en raiz y rechazo
r1 = h1.insert(5); h1.decrease_key(r1, 2); assert h1.minimum() == r1 and r1.key == 2

# 5. Corte en cascada >= 2 niveles
hc = FibonacciHeap()
A, B, C = hc.insert(1), hc.insert(10), hc.insert(20)
D, E, F, G = hc.insert(15), hc.insert(25), hc.insert(30), hc.insert(35)
hc._link(G, C); hc._link(F, C); hc._link(E, B); hc._link(C, B); hc._link(D, A); hc._link(B, A)
ini_c = hc.cascading_cuts; hc.decrease_key(E, 0); hc.decrease_key(G, 0); hc.decrease_key(F, 0)
assert (hc.cascading_cuts - ini_c) >= 2

# 6. Eliminacion de nodo interno
h_del = FibonacciHeap(); nR, nI, nH = h_del.insert(1), h_del.insert(10), h_del.insert(20)
h_del._link(nH, nI); h_del._link(nI, nR); h_del.delete(nI); h_del.validate()
\end{lstlisting}

\begin{lstlisting}[style=consolestyle, caption={Salida de consola de la batería de pruebas del Propuesto 1}]
[Prueba 1] Inserciones con repetidos y extraccion ordenada: 100% OK.
[Prueba 2] Union de monticulos vacios y no vacios: 100% OK.
[Prueba 3] Disminucion en raiz y validacion de cotas: 100% OK.
[Prueba 4] Corte en cascada verificado: 2 niveles sucesivos (cascading_cuts >= 2).
[Prueba 5] Eliminacion de nodo interno con hijos: 100% OK.
[Prueba 6] Desmarcado estricto en promocion de raiz: 100% OK.
[Prueba 7] 250 operaciones aleatorias contrastadas con heapq: 100% OK.
>>> BATERIA COMPLETA SUPERADA CON EXITO (8 INVARIANTES CERTIFICADAS).
\end{lstlisting}

\subsection{Análisis gráfico del potencial y balance amortizado}
La Figura \ref{fig:propuesto_1_potencial} ilustra la dinámica de $\Phi(H) = t(H) + 2m(H)$ durante una secuencia de inserciones seguidas de una extracción del mínimo, así como la compensación contable entre costo real y amortizado.

\begin{figure}[H]
    \centering
    \includegraphics[width=0.90\textwidth]{fig_propuesto_1_potencial.png}
    \caption{Dinámica de la función potencial $\Phi(H)$ y balance de costo real vs. amortizado}
    \label{fig:propuesto_1_potencial}
\end{figure}
